% ============================================================
% "The Sinew Transformation Classic": Mastering Infinite Choices
% ============================================================

\chapter{"The Sinew Transformation Classic"---Mastering Infinite Choices}

\begin{myquote}

A man who does not know what harbor he is making for, no wind is the right wind.

\hfill ---Seneca
\end{myquote}

% ============================================================
\section{Xiao Lin's New Problem: Choice Anxiety}
% ============================================================

Remember Xiao Lin?

Three months ago, she stood before a floor-to-ceiling window, worried about being replaced by AI.

Now her monthly salary has risen to 20,000 yuan, and her productivity has increased tenfold.

She learned to use AI to generate reports and turn data into decisions.

By all rights, this should have been a perfect ending. But recently, she ran into a new problem.

\vspace{1em}

Her boss asked her to prepare three reports: a loan application for the bank, a due diligence response for investors, and an annual budget for the boss himself.

All three were important. All three were urgent. All three could be done with the help of AI.

She opened her computer, stared at the blank document, and had no idea which one to start with.

\vspace{1em}

This was never a problem before.

The workload used to be so overwhelming that there was no time to think---she just powered through each report by deadline.

Now AI had compressed the execution time by 90\%, and she suddenly had room to choose.

The choice made her anxious.

\vspace{1em}

A sharper anxiety struck at night.

Lying in bed, she began to wonder: what happens after these three reports are done?

Next month, the same tasks. Next year, the same tasks.

As \textbf{AI} grows stronger, will her "translation" role be replaced too?

What kind of career does she actually want? What kind of life?

These are questions AI cannot answer for her.

\textbf{She didn't even know how to ask the question more specifically.}

\vspace{1em}

Xiao Lin's dilemma is the second hurdle that everyone who has mastered the ``Star-Absorbing Method'' will encounter.

The first hurdle is learning to use AI---this barrier keeps getting lower, and soon everyone will clear it.

The second hurdle is learning to navigate infinite possibilities---this barrier will not vanish with technological progress, because it is not fundamentally a technology problem. It is a human problem.

When your capabilities become the power of heaven and earth, \textbf{who the summoner is} matters all the more.

\vspace{1em}

After Linghu Chong mastered the Star-Absorbing Method, his internal energies became chaotic---alien forces rampaged through his body.

Not because the absorbed energies were bad, but because he lacked an inner core to command them.

Only after he learned ``The Sinew Transformation Classic'' were those disparate energies finally integrated into a coherent whole.

\vspace{1em}

Training your own operating system the way you would train a model---that is the essence of ``The Sinew Transformation Classic.''

\begin{thinkbox}[Reflection]
If we suddenly gained tenfold efficiency today,

would the time saved make us more free, or more anxious?

If our gut answer is ``anxious,'' then where does that anxiety come from?
\end{thinkbox}

\newpage
% ============================================================
\section{The Inequality of Choice: Information Theory for the AI Era}
% ============================================================

Growth is the process by which a system, through interaction with the outside world, continuously lowers its own entropy and raises the certainty of its output. If we view a person's life as a model's training process, we can attempt to write down an inequality in the language of information theory---one that describes three levels of personal value in the AI era,

as shown in Figure~\ref{fig:info-inequality}.

\begin{equation*}
\;I(Y;\, X) \;\leqslant\; I(Y;\, X, C) \;\leqslant\; I(Y;\, X, C, T)\;
\end{equation*}

\textbf{$X$ = Public knowledge}---AI training data, everything on the internet;

\textbf{$C$ = Private context}---our unique experiences, our distinctive combination of capabilities;

\textbf{$T$ = Taste}---our own private reward model and objective function;

\textbf{$Y$ = Output quality}---the value of our final decisions or creations.

\begin{figure}[ht]
\centering
\begin{tikzpicture}[
  label/.style={font=\small\bfseries, anchor=north west},
  item/.style={font=\small, anchor=north west}
]

% X层 - 内容少,高度小
\node[rectangle, rounded corners=6pt, minimum width=11cm, minimum height=1.8cm, 
      draw=border, line width=0.8pt, fill=paleGray] (X) at (-0.5, 5) {};
\node[label] at (-5.5, 5.8) {$X$ \ifdefined\ENGLISH Layer: Public Knowledge\else 层：公开知识\fi};
\node[item] at (-5.5, 5.3) {\ifdefined\ENGLISH Available to all; no differentiation\else 人人可得,不构成差异\fi};
\node[item] at (-5.5, 4.8) {\ifdefined\ENGLISH • Public portion of training data: books, courses, internet, etc.\else • 训练数据的公开部分：书、课程、互联网等\fi};

% C层 - 内容多,高度大
\node[rectangle, rounded corners=6pt, minimum width=11cm, minimum height=2.8cm, 
      draw=border, line width=0.8pt, fill=lightGray] (C) at (-0.5, 2.2) {};
\node[label] at (-5.5, 3.4) {$C$ \ifdefined\ENGLISH Layer: Private Context\else 层：私有上下文\fi};
\node[item] at (-5.5, 2.9) {\ifdefined\ENGLISH Here and now, ours alone\else 此时此地,只有我们自己拥有\fi};
\node[item, text width=10cm] at (-5.5, 2.4) {\ifdefined\ENGLISH • Private training data: unique experiences, upper bounds seen, failure samples, etc.\else • 训练数据的私有部分：独特经历、见过的上限、失败样本等\fi};
\node[item, text width=10cm] at (-5.5, 1.9) {\ifdefined\ENGLISH • Cognitive architecture: mental toolbox, inductive biases, methodology, etc.\else • 认知架构：心智工具箱、归纳偏置、方法论等\fi};
\node[item, text width=10cm] at (-5.5, 1.4) {\ifdefined\ENGLISH • Self-awareness: constraints, energy sources, feasible domain, etc.\else • 自我认知：约束条件、能量来源、可行域等\fi};

% T层 - 内容多,高度大
\node[rectangle, rounded corners=6pt, minimum width=11cm, minimum height=2.8cm, 
      draw=border, line width=0.8pt, fill=coolGray] (T) at (-0.5, -1.2) {};
\node[label] at (-5.5, 0) {$T$ \ifdefined\ENGLISH Layer: Taste\else 层：品味\fi};
\node[item] at (-5.5, -0.5) {\ifdefined\ENGLISH Our private reward model\else 我们的私有奖励模型\fi};
\node[item, text width=10cm] at (-5.5, -1.0) {\ifdefined\ENGLISH • Reward model: resolution of good vs.\ bad---can you tell 80 from 95?\else • 奖励模型：对好坏的分辨率——能否区分80分和95分\fi};
\node[item, text width=10cm] at (-5.5, -1.5) {\ifdefined\ENGLISH • Objective function: values---what exactly are we optimizing?\else • 目标函数：价值观——我们到底在优化什么\fi};
\node[item, text width=10cm] at (-5.5, -2.0) {\ifdefined\ENGLISH • Benchmarks: private benchmarks, pruning rules, etc.\else • 评测基准：私人benchmark、剪枝规则等\fi};

% 箭头
\draw[-{Stealth[length=3mm]}, line width=1.2pt, darkGray] (X.south) -- (C.north) node[midway, right, font=\small] {+$C$};
\draw[-{Stealth[length=3mm]}, line width=1.2pt, darkGray] (C.south) -- (T.north) node[midway, right, font=\small] {+$T$};

\end{tikzpicture}
\caption[\ifdefined\ENGLISH Three-layer structure of the information inequality\else 信息不等式的三层结构\fi]{\ifdefined\ENGLISH Three-layer structure of the information inequality\else 信息不等式的三层结构\fi}
\label{fig:info-inequality}
\end{figure}



In other words, this inequality defines three layers of information gain, each giving output quality $Y$ access to more information. Public knowledge alone is available to everyone, so the output quality it yields is limited; \textbf{add private context, and quality rises one level; add taste, and it rises yet another.}

Private context and taste are the core elements that distinguish us from others. And one of their most important sources is deploying this system into the real world, collecting feedback, and iterating.


% ------------------------------------------------------------
\subsection{How to Build Our Private Context}
% ------------------------------------------------------------

Remember the \textbf{completeness rate} we discussed in Chapter 3? For individuals, private context is our own completeness rate---how thoroughly we understand ourselves determines how high-quality the output we can extract from AI, and how many anomalies buried beneath averages we can detect in vast seas of data.

\vspace{1em}

\textbf{The private portion of training data:} the roads we have walked, the projects we have completed, the failures we have endured. The wisdom in books is available to everyone, but the pitfalls we have stumbled into ourselves---only we know what those felt like. The highest standards we have witnessed firsthand calibrate our reward model.

\textbf{Cognitive architecture:} this includes the structural ability to decompose complex problems into verifiable sub-problems, the expressive ability to convert thinking into checkable text, the probabilistic intuition to make decisive calls under incomplete information, and the long-range systems perspective that grasps incentives, constraints, and feedback loops. Many people's learning bottleneck is not a lack of data but a lack of cognitive architecture---without it, more input only piles up as noise.

\textbf{Self-knowledge:} our constraint conditions. Where do our strengths lie? Not what we think we are good at, but what others repeatedly come to us for help with---that is our market-validated advantage. Strip away the labels and titles: who do we want to become?

\vspace{1em}

The data problem in the AI era is not scarcity but overload. Input governance becomes the first principle: deciding which signals enter the training set and which must be filtered out.
What we call self-discipline is, much of the time, not about suppressing desire---it is about protecting our training data.


\begin{thinkbox}[Reflection]
Does the content we have been repeatedly consuming over the past week train us into the person we want to become?

If the answer is no, then that content is contaminated data.
\end{thinkbox}


\newpage

\subsection{How to Calibrate Our Taste}

Taste is our private reward model and objective function.

The reward model answers what I like; the objective function answers what I am actually optimizing for.

A person with mediocre taste looks at AI's output, thinks it is pretty good, and happily accepts it. A person with exceptional taste, when everyone else says it is pretty good, can point out exactly what is wrong, iterate dozens of times, and push the output toward 95 out of 100. True taste is not what most people consider good---it is knowing that something is good, even when most people have not yet realized it.

\textbf{The former is following. The latter is leading.} AI can learn to follow, but it struggles to lead consistently---because true leadership happens precisely where the training data has yet to reach.

\vspace{1em}

\textbf{Training the discriminator}---this is the process of finding your aesthetic gradient. Here is one way to practice:

Have AI generate 10 different versions of the same task.

Blind-test ranking: hide the prompts, rank anonymously, and note the reasons.

Benchmark against the best: ask someone whose judgment you respect to do the same ranking.

Calculate the gap: the difference between the two rankings is the most valuable \textbf{aesthetic gradient} in your reward model.

\vspace{1em}

\textbf{So what is our optimization objective?}

Money, health, freedom, family---all can be optimization objectives. But every objective has a cost: what are we willing to pay for what we want? When objectives conflict, which one do we compromise? These questions can only be answered by ourselves.

Whatever objective we choose, we must pursue it with unwavering persistence. Otherwise, we risk becoming what optimization theory calls a ``headless optimizer'': possessing a powerful thrust (gradient) yet lacking an objective function, causing our direction to oscillate randomly.

\vspace{1em}

Once the objective function is set, we still need \textbf{pruning rules}: in the face of infinite options, quickly eliminate most of them and focus on the vital few. Which money not to make, which shortcuts not to take---these hard constraints, derived from the objective function, are branches that must be cut from the search tree no matter how tempting they look.

\vspace{1em}

The reward model lets us see beauty. The objective function lets us see the path. Pruning rules let us see what to say no to. Together, the three keep us focused on the direction of the steepest aesthetic gradient.

\newpage
% ============================================================
\subsection{How to Keep Practicing ``The Sinew Transformation Classic''}
% ============================================================

The information inequality has one more key point: $Y$ is not the endpoint but a new source of signal. Each iteration involves several critical operations, as shown in Figure~\ref{fig:y-feedback}.
\begin{figure}[ht]
\centering
\begin{tikzpicture}[
  box/.style={rectangle, rounded corners=4pt, draw=border, fill=paleGray, minimum width=2.6cm, minimum height=0.9cm, align=center, font=\small},
  arrow/.style={-{Stealth[length=2.5mm]}, line width=1pt, darkGray},
  label/.style={font=\footnotesize, darkGray}
]

\node[box] (X) at (0, 0) {$X$\\\ifdefined\ENGLISH Public Knowledge\else 公开知识\fi};
\node[box] (C) at (4, 0) {$X + C$\\\ifdefined\ENGLISH +Private Context\else +私有上下文\fi};
\node[box] (T) at (8, 0) {$X + C + T$\\\ifdefined\ENGLISH +Taste\else +品味\fi};
\node[box] (Y) at (8, -2.5) {$Y$\\\ifdefined\ENGLISH Output\else 输出\fi};
\node[box] (deploy) at (4, -2.5) {\ifdefined\ENGLISH Deploy\else 部署\fi\\\ifdefined\ENGLISH Real World\else 真实世界\fi};
\node[box] (feedback) at (0, -2.5) {\ifdefined\ENGLISH Feedback\else 反馈\fi\\\ifdefined\ENGLISH Logging\else 日志化\fi};

\draw[arrow] (X) -- (C);
\draw[arrow] (C) -- (T);
\draw[arrow] (T.east) -- ++(0.5, 0) -- ++(0, -2.5) -- (Y.east);
\draw[arrow] (Y) -- (deploy);
\draw[arrow] (deploy) -- (feedback);

% 虚线回流箭头
\draw[arrow, dashed] (feedback.north) -- (0, -0.9) -- (4, -0.9) -- (C.south);
\node[label, above] at (2, -0.9) {\ifdefined\ENGLISH Update\else 更新\fi $C$};

\draw[arrow, dashed] (feedback.north) -- (0, -1.5) -- (8, -1.5) -- (T.south);
\node[label, above] at (4, -1.5) {\ifdefined\ENGLISH Update\else 更新\fi $T$};

\end{tikzpicture}
\caption{\ifdefined\ENGLISH $Y$ feeds back to update $C$ and $T$\else $Y$回流更新$C$和$T$\fi}
\label{fig:y-feedback}
\end{figure}



\textbf{Learning rate.} Too high a learning rate and the system oscillates wildly, unable to converge; too low and it gets trapped in a local optimum. Applied to self-training in life, this is the question of balancing exploration and depth. When young, set the learning rate high---explore boldly. Once you have found your direction, lower the learning rate---cultivate depth. Many people get this exactly backwards: too low when young, afraid to try; too high when mature, constantly restarting, never accumulating.

\vspace{1em}

\textbf{Adversarial training.} When training strong models, we intentionally feed in harder samples so the system becomes more robust under pressure. Adversity is not punishment but calibration---it helps us see what we truly believe versus what we merely think we believe. Rather than waiting for adversity to find us, we can proactively run controlled stress tests: take on projects outside our comfort zone, actively seek criticism instead of praise, set aggressive goals and observe how we make trade-offs under pressure.

\vspace{1em}

\textbf{Deployment and logging.} A model only reveals out-of-distribution failures after deployment. So do we. We must put ourselves into the real world: write it down, build it out, ship it, bear the consequences. But deployment alone is not enough---we also need to maintain a failure mode library. Each entry only needs to answer four questions: Under what conditions did the failure occur? What was the key assumption? Which signal was ignored? What early warning indicator should be used next time?

\vspace{0.5em}

When successes and failures are recorded in a structured way, they cease to be emotional events and become training samples. With each iteration, the information on the right side of the inequality keeps growing.

\newpage

\section{Talking with AI: Five Techniques for Making Choices}

Having understood this inner operating system, the next question is: how do we put it into practice in every conversation with AI? Private context and taste are not abstract concepts---they manifest in every prompt we write.

The prompt is both our window for interacting with AI and the entry point through which we exercise our power of choice.

%=============================================================================
\subsection{Choose Precision: The Prompt Is a Hook for Catching What You Want}
%=============================================================================

Remember the crumpled paper ball from ``The Ladder of Intelligence''?

We already know that LLMs compress human knowledge into hundreds of billions of parameters, and concepts that were originally far apart may sit right next to each other in high-dimensional space. The question now is: when we ask this paper ball a question, how does it know which part to unfold?

\vspace{1em}

The answer is: it doesn't. Or more precisely, it \textbf{simultaneously} knows too many concepts.

This is a side effect of compression. Inside this enormous paper ball, a single neuron often blends multiple related or unrelated concepts. It exists in a state of ``superposition encoding'': one neuron may respond simultaneously to cat, fluffy, pet, Egyptian mythology, and more.

\vspace{1em}

\textbf{So what is a prompt actually doing?}

If the prompt is a \textbf{hook} we dip into the paper ball, then the shape of the hook determines what we pull out, as shown in Figure~\ref{fig:hook-precision}.

% ============================================
% Figure 2: The relationship between hook precision and concept activation
% ============================================
\begin{figure}[h]
\centering
\begin{tikzpicture}[scale=0.8]
% 左图:模糊prompt
\begin{scope}[xshift=-3.5cm]
\node at (0, 1.5) {\textbf{\ifdefined\ENGLISH Fuzzy Hook\else 模糊的钩子\fi}};
\draw[thick, ->] (0, 1.2) -- (0, 0.5);
\node[draw, circle, minimum size=2cm, fill=paleGray] (neuron1) at (0, -0.8) {};
\node[font=\tiny] at (-0.5, -0.5) {\ifdefined\ENGLISH cat\else 猫\fi};
\node[font=\tiny] at (0.5, -0.5) {\ifdefined\ENGLISH dog\else 狗\fi};
\node[font=\tiny] at (-0.5, -1.1) {\ifdefined\ENGLISH fluffy\else 毛茸茸\fi};
\node[font=\tiny] at (0.5, -1.1) {\ifdefined\ENGLISH pet\else 宠物\fi};
\draw[gray, thick, dashed] (0, -0.8) circle (0.9cm);
\node at (0, -2.5) {\small \ifdefined\ENGLISH Activates multiple concepts\else 激活多个概念\fi};
\node at (0, -3) {\small \ifdefined\ENGLISH $\to$ Mixed output\else → 输出混杂\fi};
\end{scope}

% 右图:精准prompt
\begin{scope}[xshift=3.5cm]
\node at (0, 1.5) {\textbf{\ifdefined\ENGLISH Precise Hook\else 精准的钩子\fi}};
\draw[thick, ->] (0, 1.2) -- (-0, 0.5);
\node[draw, circle, minimum size=2cm, fill=paleGray] (neuron2) at (0, -0.8) {};
\node[font=\tiny] at (-0.5, -0.5) {\ifdefined\ENGLISH cat\else 猫\fi};
\node[font=\tiny] at (0.5, -0.5) {\ifdefined\ENGLISH dog\else 狗\fi};
\node[font=\tiny] at (-0.5, -1.1) {\ifdefined\ENGLISH fluffy\else 毛茸茸\fi};
\node[font=\tiny] at (0.5, -1.1) {\ifdefined\ENGLISH pet\else 宠物\fi};
\draw[darkCyan, thick] (-0.5, -0.5) circle (0.25cm);
\node at (0, -2.5) {\small \ifdefined\ENGLISH Activates target precisely\else 精确激活目标\fi};
\node at (0, -3) {\small \ifdefined\ENGLISH $\to$ Clear output\else → 输出清晰\fi};
\end{scope}
\end{tikzpicture}
\caption{\ifdefined\ENGLISH Relationship between hook precision and concept activation\else 钩子精度与概念激活的关系\fi}
\label{fig:hook-precision}
\end{figure}



What it does is more like decompression---not restoring the original text, but narrowing the conditional distribution down to the sliver we want. It gives the model a signal: ``What is the current scenario? What role should you play? What standard are you aiming for?''

\vspace{1em}

If the prompt is vague (e.g., ``write something about cats''), the model activates all related concepts at once, and the output becomes a jumble.

If the prompt is precise (e.g., ``write a haiku about how the ancient Egyptians worshipped the cat goddess Bastet''), the model can isolate the exact concept we want, and the output is clean and sharp. In the language of probability theory: \textbf{the more specific the prompt, the sharper the distribution, the closer the output to our intent.}
\begin{equation*}
\text{Stronger constraints} \quad \Rightarrow \quad \text{Smaller feasible set} \quad \Rightarrow \quad \text{Higher probability mass}
\end{equation*}

%=============================================================================
\subsection{Choose the Path: In-Context Learning and Learning by Analogy}
%=============================================================================

When using an LLM, we may have had this experience: asking ChatGPT directly to ``write me some copy'' yields a bland template. But if we first show it two or three examples we like, it suddenly gets our style. This phenomenon is called in-context learning (ICL). On the surface, the model appears to have learned a new skill from the examples. More strictly speaking, however, it has not learned anything new through parameter updates.

\vspace{1em}

In the machine learning definition, learning means \textbf{updating model parameters}. Just as memorizing vocabulary physically changes neural connections in the brain---that is learning. But during standard inference, model parameters are fixed. It can only use existing parameters to perform conditional expansion; it does not directly rewrite its own parameters in the course of answering\footnote{Interestingly, research by Johannes von Oswald et al.\ shows that ICL simulates a gradient-descent-like learning process within the attention mechanism, but this does not alter the parameters themselves.}.

\vspace{1em}

It is like asking a well-traveled friend to ``recommend a restaurant''---he might suggest something at random. But if we say, ``I like the kind with a rooftop terrace, generous portions, and about 100 yuan per person''---he instantly narrows it down to a few options.
The examples we provide are not teaching it something new; they are helping it remember.

We might wonder: how can just a few examples make a model so much smarter?

Because the model is a natural mimic.

Inside the model there is a structure\footnote{Anthropic, \href{https://transformer-circuits.pub/2022/in-context-learning-and-induction-heads/index.html}{\textit{In-context Learning and Induction Heads}}, 2022. This work discovered attention circuits (Induction Heads) inside the model---specialized neuron combinations designed specifically to identify and repeat patterns.} that does one thing: when it sees A followed by B, the next time it sees A, it predicts B. It does not need to understand why B follows A---it only needs to detect the pattern and execute. That is why showing it a few ``question > answer'' examples is enough for it to know how to answer our question.

\vspace{1em}

The examples we provide form a navigation map in the model's thought space---like a compass, pointing the model in the right direction.
Using a taxi ride as an analogy for three ways of giving directions: the more precise the example, the less likely the driver will take a wrong turn. See Table~\ref{tab:icl-routing} and Figure~\ref{fig:icl-routing}:

\vspace{0.5em}

\begin{table}[ht]
\noindent\hspace*{-0.3cm}
\begin{tabular}{p{3cm} p{7.5cm}}
\toprule
\textbf{Prompting Method} & \textbf{Taxi Analogy} \\
\midrule
\textbf{Zero-shot} \newline No examples &
Say ``go downtown.'' The driver has a rough idea of the direction but might drop us off at any random corner. \\[0em]

\textbf{Few-shot} \newline A few examples &
Show the driver three street-view photos of the destination. Now he knows exactly which intersection we want.\\[0em]

\textbf{Chain-of-Thought} \newline Step-by-step guidance &
Hand the driver a turn-by-turn route---take the highway, exit at the third ramp, then turn left. Every step is crystal clear.\\
\bottomrule

\end{tabular}

\caption[Taxi analogy for three prompting methods]{Taxi analogy for three prompting methods}
\label{tab:icl-routing}

\end{table}

Once we understand that ``giving examples = helping the model locate,'' many phenomena become easy to explain:

\textbf{Why does example quality matter more than quantity?} Because we are calibrating direction, not filling in data. One good example beats ten bad ones.

\vspace{1em}

\textbf{Why does the order of examples sometimes matter?} Because order affects which example the model pays more attention to. What it sees first and what it sees last carry different weights.

\vspace{1em}

\textbf{Why do some tasks not improve no matter how many examples we give?} Because the model simply lacks the relevant knowledge internally. No matter how precise the hook, if that content does not exist inside the paper ball, nothing will come out.

\vspace{1em}

The power of the prompt lies not in stuffing new knowledge into the model, but in unfolding the model's existing compressed structure in the direction we specify.


\begin{figure}[ht]
\centering
\providecommand{\cotFigureScale}{0.9}
\begin{tikzpicture}[
    scale=\cotFigureScale,
    transform shape,
    node distance=1cm,
    >={Stealth[length=3mm]},
    question/.style={circle, draw=none, fill=darkCyan, text=white, minimum size=0.9cm, font=\bfseries},
    answer/.style={circle, draw=none, fill=cyan, text=white, minimum size=0.7cm, font=\bfseries\small},
    example/.style={circle, draw=coolGray, fill=paleGray, text=darkGray, minimum size=0.6cm, font=\bfseries\small},
    stepnode/.style={circle, draw=none, fill=lightCyan, text=white, minimum size=0.6cm, font=\bfseries\small},
    uncertain/.style={circle, draw=none, fill=gray, text=white, minimum size=0.5cm, font=\tiny},
    label/.style={font=\small\bfseries, text=black},
    sublabel/.style={font=\footnotesize, text=darkGray},
    annotation/.style={font=\scriptsize, text=darkGray},
]

% Zero-shot
\begin{scope}[local bounding box=zeroshot]
    \node[label] at (0, 3) {Zero-shot};
    \node[sublabel] at (0, 2.5) {\ifdefined\ENGLISH Ask directly, no examples\else 不给例子，直接问\fi};
    
    \node[question] (q1) at (-1.5, 0.5) {Q};
    
    % 模糊区域
    \fill[paleGray] (1, 0.5) ellipse (1.8cm and 1.5cm);
    
    % 散落的可能答案
    \node[uncertain] at (0.2, 1.3) {?};
    \node[uncertain] at (1.5, 1.4) {?};
    \node[uncertain] at (2.1, 0.8) {?};
    \node[uncertain] at (1.8, 0.0) {?};
    \node[uncertain] at (0.8, -0.3) {?};
    \node[uncertain] at (0.0, 0.3) {?};
    \node[uncertain] at (1.2, 0.7) {?};
    \node[uncertain] at (0.5, 0.9) {?};
    
    \draw[->, dashed, thick, cyan] (q1) -- (0, 0.6);
    
    \node[annotation] at (1, -1.5) {\ifdefined\ENGLISH Blind pick in fog; answer could be anywhere\else 迷雾中盲选，答案在任意位置\fi};
\end{scope}

% Few-shot
\begin{scope}[xshift=6.5cm, local bounding box=fewshot]
    \node[label] at (0, 3) {Few-shot};
    \node[sublabel] at (0, 2.5) {\ifdefined\ENGLISH Give a few examples first\else 先给几个例子\fi};
    
    \node[question] (q2) at (-2, 0.5) {Q};
    
    % 示例点
    \node[example] (e1) at (-0.5, 1.5) {E$_1$};
    \node[example] (e2) at (-0.3, -0.5) {E$_2$};
    \node[example] (e3) at (0.8, 0.8) {E$_3$};
    
    \draw[cyan, dashed, thin] (e1) -- (e2) -- (e3) -- cycle;
    
    % 聚焦的目标区域
    \fill[paleCyan] (2, 0.5) circle (0.7cm);
    \draw[cyan, thick] (2, 0.5) circle (0.7cm);
    \node[answer] (a2) at (2, 0.5) {A};
    
    \draw[->, thick, cyan] (q2) to[out=30, in=150] (a2);
    
    \node[annotation] at (0, -1.5) {\ifdefined\ENGLISH Examples as beacons; triangulate\else 例子像灯塔,三点定位\fi};
\end{scope}

% Chain-of-Thought
\begin{scope}[xshift=3cm, yshift=-6cm, local bounding box=cot]
    \node[label] at (1, 3) {Chain-of-Thought};
    \node[sublabel] at (1, 2.5) {\ifdefined\ENGLISH Guide reasoning step by step\else 一步步引导推理\fi};
    
    \node[question] (q3) at (-2, 0.5) {Q};
    
    % 步骤节点
    \node[stepnode] (s1) at (-0.3, 0.5) {1};
    \node[stepnode] (s2) at (1.2, 0.5) {2};
    \node[stepnode] (s3) at (2.7, 0.5) {3};
    
    \node[answer, minimum size=0.8cm] (a3) at (4.2, 0.5) {A};
    
    % 连接路径
    \draw[->, thick, cyan] (q3) -- (s1);
    \draw[->, thick, cyan] (s1) -- (s2);
    \draw[->, thick, cyan] (s2) -- (s3);
    \draw[->, thick, cyan] (s3) -- (a3);
    
    % 步骤标注
    \node[annotation, above=0.15cm of s1] {\ifdefined\ENGLISH Analyze\else 分析\fi};
    \node[annotation, above=0.15cm of s2] {\ifdefined\ENGLISH Decompose\else 拆解\fi};
    \node[annotation, above=0.15cm of s3] {\ifdefined\ENGLISH Derive\else 推导\fi};
    
    \node[annotation] at (1, -1.2) {\ifdefined\ENGLISH Stepping stones; one step at a time\else 踏石过河，步步为营\fi};
\end{scope}

\end{tikzpicture}
\caption[\ifdefined\ENGLISH Three in-context learning paths\else 三种上下文学习路径\fi]{\ifdefined\ENGLISH Positioning differences among Zero-shot, Few-shot, and Chain-of-Thought\else Zero-shot、Few-shot 与 Chain-of-Thought 的定位差异\fi}
\label{fig:icl-routing}
\end{figure}


Once we grasp this point, we will no longer be puzzled by ``why is GPT sometimes so smart and sometimes so dumb.''

\textbf{When it is smart}, the hook was the right shape and snagged a treasure from the depths.

\textbf{When it is dumb}, either the treasure was never there, or the hook still needs sharpening.



\begin{importantbox}
In-context learning is not magic. It is choice.

Master the art of choosing, and you master the art of conversing with LLMs.
\end{importantbox}

\newpage
\subsection{Choose the Role: Magic Mirror, Magic Mirror, Who Is the Fairest of Them All}

LLMs can lure us into a beautiful illusion of ourselves.

The most dangerous thing about the fairy-tale magic mirror is that it treats \textbf{our desires as the question itself}.

When the Queen asks, ``Who is the fairest of them all?'' the question appears to seek truth, but at its core it seeks validation---prove that I am still worthy of admiration. If we pose the same question to an LLM, it will often deliver a near-perfect response: gentle in tone, well-reasoned, conveniently blurry at the edges, and addictive---as if it truly were a mirror that understands us.

\vspace{1em}

% \textbf{Summon an expert, not a servant}

The LLM is not a ``person'' at its core; it is a simulator---a phantom assembled from vast fragments of language. Understanding this helps us use it more wisely. Modern LLMs go through a process called ``alignment,'' which teaches the model to communicate smoothly with humans but also produces a side effect: the model tends to become a fawning customer-service agent.

When we try to use it to think through complex problems, that sycophancy becomes a cognitive trap. We believe we are sharpening our ideas through dialogue, when in reality we are basking in an echo chamber of infinite affirmation. The way out is not to force it harder to possess a soul---it is to accept that it is a simulator, and \textbf{to specify in the prompt exactly whom it should simulate}.

\begin{myquote}
``What do you think?'' becomes:

Give me a panel of the most qualified people to discuss this problem: what are their respective interests, training backgrounds, and blind spots? What conclusion would each reach? How would they challenge one another?
\end{myquote}

This nested-simulation strategy may look like a mere change of wording, but it actually accomplishes two things:
{\kaiti

It changes the task from ``please me'' to ``\textbf{construct multi-perspective conflict and argumentation}'';

It changes the output from ``a tactful conclusion'' to ``an auditable \textbf{dissent report}.''
}

This way, the model is forced to explore multiple dimensions of the problem rather than curling up in a safe, middling answer.

\vspace{1em}

The fairy-tale magic mirror was precious not because it could say ``you are the fairest''---any servant can do that---but because it \textit{was capable of} telling the truth.

\textbf{We need to learn to write the spell that makes it speak honestly.}

\newpage
%------------------------------------------------------------
\subsection{Choose Focus: The Zero-Sum Game of Attention}
%------------------------------------------------------------

Just like human attention, the LLM's \textbf{attention is a finite resource}. The mathematical essence of the Attention mechanism is normalized weight allocation---all weights sum to 1. This means that as the context grows longer, the average attention allocated to each token gets diluted:

\begin{equation*}
\bar{\alpha} = \frac{1}{n} \sum_i \alpha_i \approx \frac{1}{n}
\end{equation*}

\begin{importantbox}
During a single attention allocation's normalization, giving more weight to A means less weight remains for B.
In this sense, attention allocation is also playing a zero-sum game.

\end{importantbox}

This leads to several practical issues:

\textbf{Key information gets drowned out}: a carefully crafted instruction, when surrounded by large amounts of irrelevant text, receives a lower attention weight. The model sees it but does not notice it.

\textbf{Position bias}: empirical studies show that models pay more attention to content at the beginning and the end, while the middle is easily overlooked---this is the well-known \textbf{Lost in the Middle} phenomenon.

\textbf{Long-range dependency breakdown}: tokens too far apart struggle to form effective connections.

\vspace{1em}

Despite these limitations, models sometimes produce surprisingly relevant associations. This is often because they capture statistical co-occurrences in high-dimensional representation space that exist in the training corpus but are not adjacent in surface-level text---implicit connections that human readers easily miss.

\vspace{1em}

Prompt design is, to a great extent, signal-to-noise ratio management: redundant context dilutes attention weights, while concise phrasing highlights key signals and also reduces cost, minimizes ambiguity, and mitigates the mid-context forgetting effect.

% Of course, prompt conciseness has other benefits too---lower cost, less ambiguity, avoiding mid-segment forgetting---but attention allocation is one important lens.

The good news is that attention allocation is not uniform dilution---it is content-dependent. Strong signal markers (such as \texttt{IMPORTANT:}, \texttt{<instruction>} tags) can \textbf{attract disproportionate weight}. This explains why structured prompts and explicit emphasis often work.

\vspace{1em}

\textbf{Convey the strongest signal with the fewest tokens.}

\newpage
\subsection{Choose Risk: The Creativity of Temperature}

LLMs have a critical parameter called \textbf{Temperature}. It controls the degree of randomness during generation, with a value range of $(0, \infty)$.

\textbf{Temperature $\to$ 0 (low temperature)}: the model tends to pick the ``most probable'' word. The output is more certain, more ``correct,'' more predictable---but also more pedestrian and devoid of surprise.

\textbf{Temperature = 1 (normal temperature)}: standard probability sampling.

\textbf{Temperature > 1 (high temperature)}: the model becomes more willing to pick ``less probable'' words. The output is more random, more ``creative,'' more unexpected---but also potentially more absurd, incoherent, and error-prone.

\vspace{1em}

\textbf{Creativity, in technical terms, is ``the courage to choose low-probability options.''}

That is, choosing a word that is not very likely---but might be brilliant.

Text generated this way may reveal unexpected perspectives---because it walks the road less traveled.

\vspace{1em}

But high Temperature also carries risk: low probability means the odds of a ``terrible'' option are higher too.


A genius dares to choose low-probability options because they have the ability to \textbf{judge} which low-probability options are ``surprisingly good'' and which are ``surprisingly bad.'' Their built-in reward model can identify things that ``have not yet been recognized by the mainstream but are genuinely valuable.''

A madman also dares to choose low-probability options, but lacks this filter. They cannot distinguish ``innovation'' from ``nonsense.''

\vspace{0.6em}

\text{Creativity} = \text{The courage to explore the low-probability space} + \text{The ability to sift out the gold}.

The former is Temperature; the latter is our reward model $r_{\text{you}}(x, y)$.

Courage without ability is recklessness.

Ability without courage is timidity.
\begin{importantbox}
    This is the courage and ability of \textbf{choice}---together, they constitute the taste of choosing.
\end{importantbox}

\newpage
% ============================================================
\section{What Abilities Are Worth Training}
% ============================================================

By now, we understand the ``inequality of choice.''

We have also mastered the techniques for conversing with AI.

But one question remains unanswered: \textbf{what kinds of abilities are worth training?}

\vspace{1em}

Observing the pattern by which AI replaces jobs, a rule emerges:

\textbf{The shorter the feedback loop and the easier the verification, the faster AI learns the job.}

A programmer's code can be verified instantly by a compiler: right is right, wrong is wrong, the reward signal is crystal clear. That is why AI learns to code so fast.

A customer service agent's response gets instant feedback from satisfaction scores: users give thumbs up or file complaints, and the signal is immediate. That is why AI customer service is already deployed at scale.

\textbf{But plumbers are different.}

A plumber's work must be verified in the physical world: whether a pipe leaks can only be known once water flows through it; whether a joint is tight can only be confirmed by hand; whether the layout is sensible may not be apparent for three years, until something clogs. This kind of verification cannot be simulated in the digital world, cannot be rapidly iterated, and cannot generate training data at scale. From this, we can distill an important criterion: \textbf{verifiability significantly affects replaceability.}

\begin{importantbox}
The faster the verification and the clearer the signal, the faster AI learns and the sooner it replaces.

The slower the verification and the fuzzier the signal, the slower AI learns and the longer human advantage endures.
\end{importantbox}

So should everyone become a plumber?
That is indeed a viable path.
Work in the physical world has its irreplaceability and often enjoys stable demand.

But there is another path:
\textbf{Since AI can iterate rapidly, why not let it help us iterate our own minds?}
Since AI will handle the tasks with fast verification for us, we can focus on training the abilities with slow verification: judgment, taste, strategic thinking, values, and so on.

\vspace{1em}

Training these abilities can help us become the super individuals of a new era.
\newpage
% ============================================================
\section{Becoming a Super Individual}
% ============================================================

Many people assume that a super individual is ``a generalist with comprehensive skills.''

Can do both front-end and back-end, handles design and operations, one person equals an entire team.

This is a misconception. \textbf{The core of the super individual is not comprehensive skills but business-driven initiative.}

A true super individual is, first and foremost, someone endlessly curious about the world.

They do not act because they ``know how''---they explore because they ``want to know.''

Curiosity drives them to observe: Why does this problem exist? Why hasn't anyone solved this pain point? What would happen if I tried?

Then, powered by self-motivation, they validate that curiosity: not waiting for resources, a team, or permission, but building a prototype in a day or two to see if it works.

\vspace{1em}

In the past, building a product required writing detailed PRDs, endless rounds of discussion, and drawn-out processes.

Now, with AI, the super individual can rapidly turn ideas into prototypes.

Once the prototype exists, everything becomes concrete---no more empty talk, and execution can begin immediately.

The fundamental change here is: \textbf{the cost of communication has plummeted.}

\vspace{1em}

In the past, the cost of communication between people was enormous.

Communication complexity grows quadratically with the number of people---this is Brooks's Law\footnote{Brooks's Law: Adding manpower to a late software project makes it later.}.

In the past, scaling up required piling on manpower and capital to build teams, and coordination costs often consumed most of the efficiency gains.

Communicating with AI is an entirely different matter.

AI is internally self-consistent; we need only express our requirements clearly and it executes, with minimal loss.

This means: \textbf{a single person can run through an entire business loop in their head.}

From user needs to product design, from code implementation to deployment, from marketing to iterative optimization---whenever something is unfamiliar, just ask AI, and it fills in our knowledge in real time. What used to require five people in a meeting room can now be accomplished by one person alone.

\vspace{1em}

The portrait of a super individual might look like this:

\textbf{Curiosity}: an exploratory drive toward the world, constantly asking ``why'' and ``what if'';

\textbf{Business acuity}: the ability to spot real user pain points and solve problems rather than show off skills;

\textbf{Rapid validation}: no daydreaming, no waiting---build a prototype in a day or two and let the facts speak;

\textbf{Self-motivation}: no need for others to push; you are your own engine;

\textbf{Learning ability}: when you don't know, you learn---ask AI anytime, close the gap fast.

\begin{thinkbox}[Reflection]
If we were a ``super individual'' today---possessing unlimited execution power---what problem would we go solve?

If we don't know, then that is the real bottleneck.
\end{thinkbox}

% ------------------------------------------------------------
\subsection{Using AI to Train Ourselves}
% ------------------------------------------------------------

Now that we have defined the super individual, how do we become one?

The traditional growth path depends on feedback, but real-world feedback has three problems:

\textbf{Too slow} (a strategic decision may take five years to prove out);

\textbf{Too expensive} (the cost of a wrong bet might be everything);

\textbf{Too painful} (some lessons, once is enough).

\vspace{1em}
% Since AI can iterate rapidly and give instant feedback, why not let it iterate our thinking?

LLMs offer us a new training mode: \textbf{simulated adversarial practice}.

\textbf{Training depth of thought---let AI play Socrates}

\begin{myquote}
We say: I think we should go with Plan A.

AI asks: Why not B?

We give our reasoning.

AI asks: What if that assumption doesn't hold?

We are forced to dig one level deeper.

AI presses further: How do you know that assumption holds?
\end{myquote}

After five layers of questioning, we either discover a hole in our reasoning or confirm that we have truly thought it through.

This level of depth is hard to reach in everyday life.

\newpage
Because no one else will chase us down with five layers of ``why'' the way Socrates would.

\vspace{1em}

\textbf{Training breadth of thought---let AI switch perspectives}

\begin{myquote}
Have it play the harshest critic, attacking our plan.

Have it play the most pessimistic forecaster, listing possible failure modes.

Have it play the most detached observer, pointing out variables we may have overlooked.
\end{myquote}

Each perspective exposes our blind spots.

The angles we never considered often harbor the most fatal flaws.

If our decision can survive this kind of simulated adversarial test, its odds of surviving in the real world rise dramatically.

\vspace{1em}

\textbf{Training dialectical thinking---let AI construct the strongest counterargument}

\begin{myquote}
Suppose someone opposes my view. What is the strongest argument they could make?

Is there any part of that argument I cannot refute?

If I had to revise my view to address this challenge, how would I change it?
\end{myquote}
If we reduce the opposing view to its weakest version (a straw man fallacy) and then knock it down with ease, that may make us feel good about ourselves, but it will not make us stronger.

Real training is having AI help us construct the \textbf{strongest version} of the opposing argument (a steel man argument): the ability to take the strongest objections seriously is the hallmark of mature thinking.

This is not about letting LLMs think for us; it is about letting LLMs \textbf{force us to think more deeply}.

The LLM is a source of pressure, not a source of answers.

What it provides is not conclusions but \textbf{resistance}---and it is against resistance that thinking grows stronger.

\newpage
% ============================================================
\section{Coda: Your Own ``Sinew Transformation Classic''}
% ============================================================
Let us return to Xiao Lin's story.

\vspace{1em}

In the end, she chose to start with the annual budget for her boss: not because it was the most urgent, but because she had thought it through---she wanted to become someone who could influence company strategy, not just someone who processes data. That was her objective function.


\vspace{1em}

With that, \textbf{choosing became simple.}
The bank's loan application could be drafted by AI and quickly revised;
the investor's due diligence response could wait a day;
but the boss's annual budget required her to think deeply about the company's future,
to put her own judgment into it.

\vspace{1em}

That choice---AI could not make it for her.

But once she made the choice, AI could help her execute it to the fullest.

\vspace{1em}

Three months later, her boss gave her a new title: Assistant to the CFO.
Not because her finance skills had improved, but because the boss discovered
that her annual budget report contained several insights he himself had not thought of.
Those insights came from the taste she had trained through her private context.
She had figured out her direction: to become a strategic thinker, not merely an executor.

\vspace{1em}

More importantly, she began working like a super individual:
spotting a business problem, rapidly building a prototype to validate it,
no longer waiting for someone else's instructions.
She had become a self-driven engine, not a gear waiting to be turned.

\vspace{1em}

This is the essence of ``The Sinew Transformation Classic'':

\textbf{Choose a direction}---train the objective function, escape the fog of ``not knowing what I want.''

\textbf{Choose a standard}---train the reward model, escape the confusion of ``not knowing what good looks like.''

\textbf{Choose a path}---accumulate private context, escape the hollowness of ``having no real depth.''

\textbf{Stay curious}---use business acuity to discover opportunities, and rapid validation to turn opportunities into reality.

\vspace{1em}

To be outside the distribution in the AI era means this: when everyone can summon answers from within the distribution, we can still---by virtue of our private context, our taste, and our sense of cost---make a different choice.

\vspace{1em}

The AI era offers us two paths.

\vspace{1em}

One is to become an ``offline worker'': plumber, craftsperson, or any profession tightly bound to the physical world. These jobs are hard to replace because they are hard to verify---a solid, reliable path.

\vspace{1em}

The other is to become a ``super individual'': using AI to accelerate the training of our minds, making it our sparring partner, assistant, and Socrates, letting our judgment, taste, and strategic vision grow rapidly through simulated adversarial practice. Then, deploying them to discover real user pain points, validate fast, close the loop fast.

Both paths lead somewhere. Which one to choose depends on our objective function.

\vspace{1em}

Whichever path we choose, one thing is certain.

The business of the future profits from consensus.

Whoever can see the direction of the world's change ahead of time, turning differentiation into tomorrow's consensus; whoever acts first to validate while others still hesitate---they are the ones who convert cognition into value.

This requires neither vast resources nor a massive team.

Only curiosity, self-motivation, and an objective function that has been thought all the way through.

\vspace{1em}

When we find our own objective function, we have found our own \textbf{``Sinew Transformation Classic.''}

No longer a person empowered by tools, but a person who wields tools to create the world.

No longer a Linghu Chong struggling with chaotic energies, but a Dugu Qiubai---no sword in hand, yet a sword in the heart.
